\documentclass[10pt,a4paper]{amsart}
\usepackage[margin=19mm]{geometry}
\usepackage{amsmath,amssymb,mathtools}
\usepackage[scr=rsfs,cal=euler]{mathalfa}
\usepackage{xfrac}
\usepackage[unicode,hidelinks,bookmarks=false]{hyperref}
\newcommand{\ie}{i.e.\ }
\newcommand{\cf}{cf.\ }
\newcommand{\resp}{respectively\ }
\newcommand{\Subsection}{\subsection}
\providecommand{\nobreakdash}{\nobreak\hbox{-}}
\setlength{\emergencystretch}{2em}
\multlinegap0pt
\usepackage{amssymb}
\usepackage{mathtools}
\newtheorem{theorem}{Theorem}
\newcommand{\R}{\mathbb{R}}
\newcommand{\Tr}{\operatorname{Tr}}
\title{Information Loss under Gauge Reduction\\ of Discrete Electromagnetic Fields}
\author{Jean-Pierre Magnot}
\date{Source: arXiv:2609.05473v1 (CC BY 4.0)}
\begin{document}
\maketitle

\noindent\emph{Source and licence.} Information Loss under Gauge Reduction\\ of Discrete Electromagnetic Fields. Version arXiv:2609.05473v1 (CC BY 4.0), distributed by arXiv under the Creative Commons Attribution 4.0 International licence, http://creativecommons.org/licenses/by/4.0/, as recorded in the arXiv licence metadata for this exact version and shown on the abstract page https://arxiv.org/abs/2609.05473v1. Full licence notice: \texttt{licenses/arxiv-2609.05473-CC-BY-4.0.txt}.\par\smallskip

\noindent\textit{Editorial scope.} Counted unit: the article's theorem thm:explicit (source0.tex lines 614--671). The article's complete original proof of it is delivered with it (source0.tex lines 673--758, one \texttt{proof} environment of 86 lines). No mathematics is written, repaired or re-derived: the statement and the proof are the article's own lines, in the article's own notation, and no retained line is substituted or re-flowed.  Retained uncounted as the source-local context the proof needs: lines 596--603; lines 605--612.  Nothing was omitted from any retained span, so the package carries no omission record.  No retained line contains a citation, so the wrapper carries no bibliography and no bibliography entry.  No retained line contains a figure environment, an image-inclusion command or drawing code, so no image is delivered and the package declares no asset dependency.  Editorial formatting only, in addition: an AMS \texttt{amsart} preamble that restates the article's own declarations and macros used by the retained lines, namely the environments \texttt{theorem} on separate counters, together with the standard packages those lines need.  The retained environments print in this document as Theorem 1 in order of appearance, whereas the article prints the same items with the numbering of its own full document.  The complete native source of the article as distributed in the arXiv e-print for this exact version is bundled unchanged as \texttt{sources/209439/source0.tex}.
\begin{equation}
r_\mu(x)
=
m_{\mu,\mathrm g}
+
R_\mu(x-m_{\mu,\mathrm p})
\label{eq:rmu}
\end{equation}

\begin{equation}
r_\nu(x)
=
m_{\nu,\mathrm g}
+
R_\nu(x-m_{\nu,\mathrm p}).
\label{eq:rnu}
\end{equation}

\begin{theorem}
\label{thm:explicit}
Set
\begin{equation}
\Delta R
=
R_\mu-R_\nu
\label{eq:delta-R}
\end{equation}
and
\begin{equation}
d
=
m_{\mu,\mathrm g}
-
m_{\nu,\mathrm g}
-
R_\nu
\left(
m_{\mu,\mathrm p}
-
m_{\nu,\mathrm p}
\right).
\label{eq:d}
\end{equation}

Then
\begin{align}
\mathcal I_{\mathrm g}(\mu,\nu)
=
\frac12
\Bigg[
&
\Tr(S_\nu^{-1}S_\mu)
-
q
+
\log
\frac{
\det S_\nu
}{
\det S_\mu
}
\nonumber\\
&+
d^\top S_\nu^{-1}d
\nonumber\\
&+
\Tr\left(
S_\nu^{-1}
\Delta R
A_\mu
\Delta R^\top
\right)
\Bigg].
\label{eq:explicit-gauge-info}
\end{align}
\end{theorem}

\begin{proof}
For nondegenerate Gaussian measures on \(\R^q\),
\begin{align}
&D\left(
\mathcal N(r_1,S_1)
\middle\|
\mathcal N(r_2,S_2)
\right)
\nonumber\\
&\quad=
\frac12
\left[
\Tr(S_2^{-1}S_1)
-
q
+
\log
\frac{
\det S_2
}{
\det S_1
}
+
(r_1-r_2)^\top
S_2^{-1}
(r_1-r_2)
\right].
\label{eq:gaussian-relative}
\end{align}

By \eqref{eq:rmu} and \eqref{eq:rnu},
\begin{equation}
r_\mu(x)-r_\nu(x)
=
d+\Delta R(x-m_{\mu,\mathrm p}).
\label{eq:r-difference}
\end{equation}

Under \(\mu_{\mathrm p}\),
\begin{equation}
\mathbb E_{\mu_{\mathrm p}}
\left[
x-m_{\mu,\mathrm p}
\right]
=
0
\label{eq:centered}
\end{equation}
and
\begin{equation}
\mathbb E_{\mu_{\mathrm p}}
\left[
(x-m_{\mu,\mathrm p})
(x-m_{\mu,\mathrm p})^\top
\right]
=
A_\mu.
\label{eq:covariance}
\end{equation}

Therefore,
\begin{align}
&
\mathbb E_{\mu_{\mathrm p}}
\left[
(r_\mu(x)-r_\nu(x))^\top
S_\nu^{-1}
(r_\mu(x)-r_\nu(x))
\right]
\nonumber\\
&\qquad=
d^\top S_\nu^{-1}d
+
\Tr\left(
S_\nu^{-1}
\Delta R
A_\mu
\Delta R^\top
\right).
\label{eq:expectation-quadratic}
\end{align}

Integrating \eqref{eq:gaussian-relative} with respect to
\(\mu_{\mathrm p}\) gives
\eqref{eq:explicit-gauge-info}.
\end{proof}

\end{document}
